\begin{remark} Lemma \ref{lem:bourgain_lemma_2} is not true without the finite index assumption.
As a counterexample, we let $\phi(x) = x, \psi(x) = 2x$ and $G=\F_2 ^k$ for some large $k$. Let $n=|G| = 2^k$ and $\{\gamma_i: i \in [n]\}$ be the set of characters of $G$. Let $a_1, \ldots, a_n$ and $b_1, \ldots, b_n$ be nonnegative real numbers. The exact values of $a_i, b_i$ will be chosen later. For each $i \in [n]$, define 
\begin{itemize}
\item $\widehat{f_1}(\gamma_i) = \widehat{f_3}(\gamma_i) =1$,
\item $\widehat{f_2}(\gamma_i) = a_i$,
\item $\widehat{K}(\gamma_i) = b_i$.
\end{itemize}
It follows that $f_1(x) = f_3(x) = n \cdot 1_{x=0}$.
Then \eqref{eq:k} says that
\[
n(a_1 b_1 + \cdots + a_n b_n)^2 \ll (a_1^2 + \cdots + a_n^2) (b_1 + \cdots + b_n)^2.
\]
This is false by taking $a_1=b_1=1$ and $a_i = b_i =0$ for $i \neq 1$.
\end{remark}

While the previous lemma involves the configuration $x, x + \phi(y), x + \psi(y)$, the next one is concerned with $x, x + \phi(y)$ and $\psi(y)$. Its proof is almost identical and so we only highlight the differences. 

\begin{lemma} \label{lem:counting2}
 Let $\phi, \psi: G \rightarrow G$ be continuous homomorphisms such that $\phi(G), \psi(G)$ have finite index in $G$. Then for $f_1, f_2, f_3 \in L^\infty(G)$, we have
\begin{equation} 
\left| \iint_{G^2} f_1(x) f_2(x+ \phi(y))f_3(\psi(y)) \, d\mu(x) d\mu(y) \right| \ll \| \widehat{f_1} \|_{\infty} \|f_2 \|_2 \| f_3\|_2 
\end{equation}
where the implicit constant depends only on the indexes of $\phi(G)$ and $\psi(G)$ in $G$.
\end{lemma}
\begin{proof}
%\Hnote{Similar (and simpler) than the previous lemma, using Cauchy-Schwarz and Plancherel, etc.}

%The proof of this lemma is similar to the one of \cref{lem:bourgain_lemma_2} and so we just emphasize the places that they are different. 
Similar to the proof of \cref{lem:bourgain_lemma_2}, without loss of generality, we can assume $f_1, f_2, f_3$ are equal to their Fourier series. For $x \in G$, write $g(x) = \int_{G} f_2(x+ \phi(y) )f_3(\psi(y) ) \, d\mu(y)$ and then by Parseval's theorem,
\begin{equation*} %\label{eq:k2}
\left| \int_G f_1(x) g(x) \, d \mu(x) \right|= \left| \sum_{\gamma \in \Gamma} \widehat{f_1}(\gamma) \widehat{g} (\overline{\gamma}) \right| \leq \| \widehat{f_1} \|_{\infty} \cdot \| \widehat{g} \|_{1}.
\end{equation*} 
Moreover, we also have
\begin{eqnarray*}
g(x) %&=& \int_{G} f_2(x+ \phi(y)) f_3(\psi(y)) \, d\mu(y) \nonumber \\
&=& \int_G \left( \sum_{\gamma_2, \gamma_3 \in \Gamma} \widehat{f_2}(\gamma_2) \gamma_2(x+ \phi(y)) \widehat{f_3}(\gamma_3) \gamma_3(\psi(y)) \right) \, d\mu(y) \nonumber \\
%&=& \int_G \left( \sum_{\gamma_2, \gamma_3 \in \Gamma} \widehat{f_2}(\gamma_2) \widehat{f_3}(\gamma_3) \gamma_2(x) (\gamma_2 \circ \phi +  \gamma_3 \circ \psi)(y)\right) \, d\mu(y) \nonumber \\
&=&  \sum_{\substack{ \gamma_2, \gamma_3 \in \Gamma,\\ \gamma_2 \circ \phi + \gamma_3 \circ \psi = 0}} \widehat{f_2}(\gamma_2) \widehat{f_3}(\gamma_3) \gamma_2(x). %\label{eq:k3}
\end{eqnarray*}
As a consequence,
\begin{equation*} %\label{eq:k4}
\widehat{g}(\gamma) = \sum_{\substack{ \gamma_2, \gamma_3 \in \Gamma,\\ \gamma_2 \circ \phi + \gamma_3 \circ \psi = 0, \\ \gamma_2 = \gamma}} 
\widehat{f_2}(\gamma_2) \widehat{f_3}(\gamma_3) =
\widehat{f_2}(\gamma) \sum_{\substack{\gamma_3 \in \Gamma, \\ \gamma \circ \phi + \gamma_3 \circ \psi = 0}}  \widehat{f_3}(\gamma_3)
\end{equation*}
and so
\begin{equation*} %\label{eq:k5}
\| \widehat{g} \|_{1} \leq \sum_{\substack{\gamma_2, \gamma_3 \in \Gamma, \\ \gamma_2 \circ \phi + \gamma_3 \circ \psi = 0}} \left| \widehat{f_2}(\gamma_2) \right| \left| \widehat{f_3}(\gamma_3) \right|.
%&=& \sum_{\gamma_2, \gamma_3 \in \Gamma} |\widehat{f_2}(\gamma_2) | \cdot | \widehat{f_3}(\gamma_3) | \cdot 
%| \widehat{K}(-\gamma_2 - 2 \gamma_3) |
\end{equation*}
% Therefore, it suffices to show that for each $\gamma_0 \in \Gamma$, we have
% \[
% \sum_{\substack{ \gamma_2, \gamma_3 \in \Gamma,\\ \gamma_2 \circ \phi + \gamma_3 \circ \psi + \gamma_0=0}} |\widehat{f_2}(\gamma_2) | \cdot | \widehat{f_3}(\gamma_3) | \ll \|f_2 \|_2 \cdot \| f_3 \|_2.
% \]
On the other hand, we have
\begin{eqnarray}
\sum_{\substack{\gamma_2, \gamma_3 \in \Gamma, \\ \gamma_2 \circ \phi + \gamma_3 \circ \psi = 0}} \left| \widehat{f_2}(\gamma_2) \right| \left| \widehat{f_3}(\gamma_3) \right| &=& \sum_{\gamma_2 \in \Gamma} \left( \left| \widehat{f_2}(\gamma_2) \right| \sum_{\gamma_3 \in \Gamma, \atop{\gamma_3 \circ \psi = - \gamma_2 \circ \phi}} \left| \widehat{f_3}(\gamma_3) \right|\right) \nonumber\\
&\leq& \left( \sum_{\gamma_2 \in \Gamma} \left| \widehat{f_2}(\gamma_2) \right|^2 \right)^{1/2} \left( \sum_{\gamma_2 \in \Gamma} \left( \sum_{\gamma_3 \in \Gamma, \atop{\gamma_3 \circ \psi = - \gamma_2 \circ \phi}} \left| \widehat{f_3}(\gamma_3) \right| \right)^2 \right)^{1/2} \nonumber \\
&\ll& \lVert f_2 \rVert_2 \left( \sum_{\gamma_2 \in \Gamma} \sum_{\gamma_3 \in \Gamma,\atop{\gamma_3 \circ \psi = - \gamma_2 \circ \phi}} \left|  \widehat{f_3}(\gamma_3) \right|^2 \right)^{1/2} \label{eq:new_index1} \\
&\ll& \lVert f_2 \rVert_2 \left( \sum_{\gamma_3 \in \Gamma} \left| \widehat{f_3}(\gamma_3) \right|^2 \right)^{1/2} \label{eq:new_index2} \\
&=& \lVert f_2 \rVert_2 \lVert f_3 \rVert_2. \nonumber
\end{eqnarray}
In \eqref{eq:new_index1}, we use the fact that for each $\xi \in \Gamma$, there are $\leq [G: \psi(G)]$ values of $\gamma_3$ such that $\gamma_3 \circ \psi = \xi$ while \eqref{eq:new_index2} follows from the fact that there are $\leq [G: \phi(G)]$ values of $\gamma_2$ such that $\gamma_2 \circ \phi = \xi$.
\end{proof}
